% titel: Blatt 5, A4: Intervallhalbierungsverfahren
% skript: Kap. 3 (Zwischenwertsatz, numerische Nullstellensuche)

% Teil I: Werkzeuge, kein Lösungsweg.

Eine Nullstelle eines Polynoms $p$ durch fortgesetztes Halbieren eines
Intervalls approximieren.

\begin{werkzeugkasten}[Warum funktioniert Halbieren? (Zwischenwertsatz)]
Eine stetige Funktion, die an $a$ negativ und an $b$ positiv ist, muss
\emph{irgendwo zwischen} $a$ und $b$ durch $0$ gehen -- sie kann nicht
\emph{von negativ nach positiv springen}, ohne auf dem Weg $0$ zu berühren.
Das ist der \textbf{Zwischenwertsatz}, und er garantiert hier: Mit
$p(a)\cdot p(b)<0$ (genau ein Vorzeichenwechsel am Einschlusspunkt) gibt es
\emph{mindestens eine} Nullstelle in $[a,b]$.

Das Halbieren nutzt diese Idee direkt aus: Im Mittelpunkt $m=\frac{a+b}{2}$
muss $p(m)$ entweder dasselbe Vorzeichen wie $p(a)$ haben (dann liegt die
Nullstelle in $[m,b]$) oder dasselbe wie $p(b)$ (dann in $[a,m]$). Das
Prüfzeichen ist das \emph{Produkt} $p(a)\cdot p(m)<0$. In jedem Fall wird
das Intervall vertauscht mit einem halb so langen, das die Nullstelle weiter
einschliesst.

\textit{Skript:} Kap. 3 (Zwischenwertsatz).
\end{werkzeugkasten}

\begin{werkzeugkasten}[Wie schnell wird es genau?]
Jeder Schritt halbiert die Länge. Nach $k$ Schritten ist das einschliessende
Intervall $\dfrac{b-a}{2^k}$ lang; die Nullstelle $x^*$ ist bis auf die
halbe Intervallbreite bestimmt. Der entscheidende Punkt: Die Genauigkeit
wächst \emph{exponentiell mit der Zahl der Schritte}. Schon
\begin{itemize}
  \item $k=10$ Schritte: Breite $\tfrac{2}{2^{10}}\approx2\cdot10^{-3}$
        (etwa zwei Nachkommastellen),
  \item $k=20$ Schritte: Breite $\tfrac{2}{2^{20}}\approx2\cdot10^{-6}$
        (etwa fünf Nachkommastellen).
\end{itemize}
Jeweils 3 weitere Schritte liefern eine Nachkommastelle mehr. Das
Verfahren braucht keinen Ableitungsblick, nur Stetigkeit und einen
Vorzeichenwechsel -- deshalb ist es sehr robust.

\textit{Skript:} Kap. 3 (numerische Verfahren).
\end{werkzeugkasten}

\begin{achtung}[Typische Fallen]
\begin{itemize}
  \item \textbf{Vorbedingung prüfen:} $p$ stetig \emph{und}
        $p(a)\cdot p(b)<0$ als erster Schritt. Für $p(x)=x^5-2x-1$:
        $p(0)=-1<0$, $p(2)=27>0$ -- erfüllt.
  \item \textbf{Das Vorzeichen des Produkts} $p(a)\cdot p(m)$ entscheidet,
        nicht das von $p(m)$ allein (ein einziges Vorzeichen reicht nicht,
        man braucht den Vergleich mit dem linken Rand).
  \item \textbf{Sauber protokollieren:} linke/rechte Grenze und $p(m)$ bei
        jedem Schritt festhalten, dann bleibt nachvollziehbar, welches
        Teilintervall gewählt wurde.
  \item \textbf{Beim Code:} die Schleife verwaltet die Grenzen; der Zähler
        und der Produkt-Test $p(a)\cdot p(m)<0$ müssen stimmen, sonst
        konvergiert die Folge in ein falsches Intervall.
\end{itemize}
\end{achtung}
